\subsection{Generic Brownian contacts}

Throughout this subsection, $Z=(L,R)$ is a two-sided Brownian motion with
unit coordinate variances and correlation $\rho\in(0,1)$.  A common constant
variance factor changes none of the assertions.  Set
\[
 \theta=\arccos(-\rho)\in(\pi/2,\pi),\qquad
 d_2=\frac{\pi}{2\theta}>\frac12.
\]
For a continuous scalar function $f$, a \emph{horizontal chord} is a pair
$s<t$ such that
\[
 f(s)=f(t)=\min_{[s,t]}f.
\]
A \emph{branch triple} consists of three distinct times $a<b<c$ for which
$f(a)=f(b)=f(c)=\min_{[a,c]}f$.  Thus its middle time is a local minimum.
All the following assertions concern nonzero chord spans and strictly
positive witnessing time intervals.

\begin{lemma}[Images of Brownian record times]\label{arm:record-height-images}
Fix a compact time interval $K$ and $\delta>0$.  The surrounding Brownian
path is retained far enough to include every arm of length $\delta$ used
below.  For signs $\sigma,\tau\in\{-1,1\}$, let
\[
 T_{\sigma,\tau}=
 \left\{t\in K:
 \begin{array}{l}
 L(t+\sigma u)\ge L(t),\\
 R(t+\tau u)\ge R(t)
 \end{array}
 \text{ for every }0\le u\le\delta\right\}.
\]
If $\sigma=\tau$, the upper box dimensions of $Z(T_{\sigma,\tau})$
and of either coordinate projection are at most $2(1-d_2)<1$.
If $\sigma\ne\tau$, these upper box dimensions are zero.
Furthermore, almost surely there is no time $t\in K$ at which $L$ has
lower-wall arms of length $\delta$ on both sides and $R$ has such an arm
on either one of the two sides.  The same statement holds with $L,R$
interchanged.
\end{lemma}

\begin{proof}
We first record the required survival bounds.  For fixed $v>0$ and
$0<r\le1$, the reflection principle gives
\[
 \mathbb P\left[\inf_{0\le u\le v}L(u)\ge-r\right]\le C_vr.
\]
Whitening the covariance sends the two lower walls to a wedge of opening
$\theta$.  Lemma~\ref{arm:wedge-survival} and Brownian scaling give
\begin{equation}\label{arm:generic-wedge-survival}
 \mathbb P\left[\inf_{0\le u\le v}L(u)\ge-r,
                  \inf_{0\le u\le v}R(u)\ge-r\right]
 \le C_v r^{2d_2}.
\end{equation}
Only this elementary Brownian survival estimate is used here; no lattice
or passage-kernel estimate is involved.

Partition $K$ into intervals $I=[s,s+h]$ of length $h=2^{-n}$, allowing
two extra intervals at its ends.  Fix $\varepsilon\in(0,1/2)$ and put
$r=h^{1/2-\varepsilon}$.  Almost surely, for all sufficiently large $n$,
both coordinate oscillations on every such interval are at most $r$.
For example this follows from Brownian H\"older continuity with exponent
$1/2-\varepsilon/2$.  Suppose such an interval contains a record time
$t$.  If a required arm points forward, then its constraint implies
\[
 \inf_{0\le u\le\delta/2}
       \{Z_i(s+h+u)-Z_i(s+h)\}\ge-r;
\]
if it points backward, the analogous constraint is
\[
 \inf_{0\le u\le\delta/2}
       \{Z_i(s-u)-Z_i(s)\}\ge-r.
\]
These implications hold when $h\le\delta/2$.  They use deterministic
walls $-r$ relative to the respective cell endpoints, so their probability
bounds do not condition on a selected record time.

For two arms on the same side, Equation
\eqref{arm:generic-wedge-survival} bounds this necessary exterior event
by $Cr^{2d_2}$.  For arms on opposite sides the two events use disjoint
time increments, and the reflection principle bounds their product by
$Cr^2$.  In the last assertion of the Lemma, the necessary exterior event
uses two coordinate arms on one side and one coordinate arm on the other;
its probability is at most $Cr^{2d_2+1}$.

Here is the resulting covering calculation, including the dimension
needed below.  If the survival power is $\alpha$, the number $N_h$ of
cells satisfying the necessary exterior event obeys
\begin{equation}\label{arm:record-cell-count}
 \mathbb E N_h\le C h^{-1}r^\alpha
          =C h^{-1+(1/2-\varepsilon)\alpha}.
\end{equation}
For $\alpha\le2$ and every $\zeta>0$, Markov's inequality and
Borel--Cantelli imply, along the dyadic sequence,
\[
 N_h\le h^{-1+(1/2-\varepsilon)\alpha-\zeta}
 \quad\text{eventually almost surely}.
\]
Each occupied cell contributes a set of diameter at most $\sqrt2r$ to
the $Z$-image and a set of diameter at most $r$ to a coordinate image.
Consequently their upper box dimensions are at most
\[
 \frac{1-(1/2-\varepsilon)\alpha+\zeta}
      {1/2-\varepsilon}.
\]
Letting $\varepsilon,\zeta$ decrease to zero through countable sequences
gives $2-\alpha$.  Taking $\alpha=2d_2$ or $\alpha=2$ proves the two
image assertions.  If $\alpha=2d_2+1>2$, choose $\varepsilon$ so small
that $(1/2-\varepsilon)\alpha>1$.  The right side of Equation
\eqref{arm:record-cell-count} is summable over dyadic $h$, so eventually
there are no candidate cells.  This proves the nonexistence assertion.
\end{proof}

\begin{lemma}[Generic horizontal contacts]\label{arm:generic-chords}
Almost surely, simultaneously on all compact time intervals:
\begin{enumerate}[label=\textup{(\roman*)}]
 \item The local minima of each coordinate are strict, countable, and
 have pairwise distinct values.  Every scalar horizontal equivalence
 class has at most three distinct times.
 \item No nontrivial horizontal chord is a chord of both $L$ and $R$.
 \item No member of a branch triple of one coordinate is an endpoint
 of a nontrivial horizontal chord of the other coordinate.
\end{enumerate}
The assertions also hold under any law having joint local absolute
continuity with these Brownian paths, including their relative coordinate
heights on separated windows.  In particular they hold in strict interior
windows of a nondegenerate correlated Brownian quadrant excursion.
\end{lemma}

\begin{proof}
We shall use the following precise source of the height densities.
Take deterministic intervals $I=[p,q]$ and $J=[r,v]$ with $q<r$.
The two paths
\[
 (Z(t)-Z(q))_{t\in I},\qquad (Z(t)-Z(r))_{t\in J}
\]
and the increment $G=Z(r)-Z(q)$ are independent.  The last has a
nonsingular Gaussian density on $\mathbb R^2$.  Its coordinate projections
have Gaussian densities on $\mathbb R$.  Thus, conditional on the two
displayed paths, a prescribed Lebesgue-null set of relative heights is
missed almost surely.  Conditioning is only on increments inside the
two separated windows; it is never on a putative contact or chord.

For completeness this also proves the scalar minimum facts.  The minimum
values on two disjoint compact rational intervals have probability zero
to agree: express the left minimum relative to its right endpoint and
the right minimum relative to its left endpoint, and condition on both
window paths.  Equality would prescribe one value for the independent
scalar Gaussian increment in the gap.  Taking a countable union shows
that no such pair of rational intervals has equal minima.  Two distinct
minimizers on a compact rational interval could be enclosed in disjoint
compact rational subintervals of that interval with the same minimum,
so each rational compact interval
has a unique minimizer.  Every local minimum is the minimizer on one
of these intervals.  This makes it strict, makes the set countable, and
also shows that distinct local minima have distinct values.  If a scalar
horizontal class contained four ordered visits, its middle two visits
would be distinct local minima at the same height.  This proves (i).

Suppose that $s<t$ were a chord in both coordinates.  Choose disjoint
rational windows $I,J$ containing $s,t$ in their interiors, and a
positive rational $\delta$ small enough that the forward arms from $s$
and the backward arms from $t$ of length $\delta$ lie in these windows.
Such choices are countable and cover every proposed nonzero chord.
Define compact sets of relative heights using all these local joint
record times in the respective windows:
\[
 A=\{Z(u)-Z(q):u\in I\text{ has both forward arms of length }\delta\},
\]
\[
 C=\{Z(w)-Z(r):w\in J\text{ has both backward arms of length }\delta\}.
\]
Here the record times are restricted to closed subintervals leaving
the indicated arm margins; countably many choices of these subintervals
suffice.  By Lemma~\ref{arm:record-height-images}, both sets have upper
box dimension at most $2(1-d_2)<1$.  Hence their difference set $A-C$
has two-dimensional Lebesgue measure zero.  Explicitly, for any
$a\in(2(1-d_2),1)$ each set has an $\epsilon$-cover with
$O(\epsilon^{-a})$ balls; pairwise differences cover $A-C$ with
$O(\epsilon^{-2a})$ balls, whose total area tends to zero.
The endpoint equality $Z(s)=Z(t)$ would imply $G\in A-C$.
The conditional Gaussian density therefore excludes the proposed
common chord, proving (ii).

Now let $a<b<c$ be an $L$-branch triple, and suppose an $R$-chord has
an endpoint $t\in\{a,b,c\}$ and another endpoint $u\ne t$.
If $t=b$, the $L$-coordinate has lower-wall arms on both sides of $b$,
and the $R$-chord supplies a lower-wall arm on the side pointing toward
$u$.  Truncate all three arms to a common positive rational length.
The final assertion of Lemma~\ref{arm:record-height-images} excludes
this case, since $d_2+1/2>1$.

If $t=a$ or $t=c$, retain a small window around $t$ disjoint from a
small window around $b$.  At $t$ the $L$-arm points toward $b$, and the
$R$-arm points toward $u$.  Truncate both to a common positive rational
length so that they lie in the outer-visit window.  They can point in
the same direction or in opposite directions; both possibilities are
included.  In relative $L$-heights, let $A$ be the image of all local
joint record times of this specified orientation in the outer-visit
window.  Lemma~\ref{arm:record-height-images} makes $A$ a
one-dimensional Lebesgue-null set: its upper box dimension is at most
$2(1-d_2)<1$ in the first case and zero in the second.

In the middle-visit window let $C$ be the countable set of minimum
values on its compact rational subintervals, measured relative to the
window's reference endpoint.  Since $b$ is a local minimum, its
relative $L$-height belongs to $C$.  Use the right endpoint as reference
for the earlier window and the left endpoint for the later window.
The equality $L(t)=L(b)$ would force the scalar Gaussian increment
between the windows to belong to either $A-C$ or $C-A$.
Both are Lebesgue-null, being countable unions of translates of a
null set.  Conditional on the two window paths the gap increment
has a density, so the equality has probability zero.  Countably many
rational windows, arm lengths, and orientations cover every outer
contact.  This proves (iii); interchanging the two coordinates gives
the other color.  The argument allows $u$ to equal another member of
the branch triple, as well as to lie between or outside its visits.

Every exclusion just proved is a null event specified in one window
or in finitely many separated windows together with their relative
heights.  It therefore transfers under the stated joint local absolute
continuity.  On strict interior windows a Brownian quadrant excursion
has this absolute continuity by its killed transition densities and
the Markov bridge construction, first restricting to compact positive
endpoint ranges and then exhausting them.  The excursion's common
initial/final endpoint is outside the assertion.
\end{proof}

\begin{remark}[Scope of the generic-contact assertion]
\label{arm:generic-chords-scope}
For an $L$-branch triple $a<b<c$, assertion (i) also gives
$L>L(b)$ on $(a,b)\cup(b,c)$ and lower values arbitrarily close just
outside $a$ and $c$.  Indeed an additional interior visit would give
four visits, while failure of an outside lower dip would make $a$ or
$c$ a second local minimum at the height of $b$.  Adjacent contacts
are consequently approached by ordinary chords at heights decreasing
to $L(b)$ from above, and the outer contact is approached by chords
at heights increasing to $L(b)$ from below.  This is a scalar consequence
of continuity and these strict inequalities.

The proof does not exclude a time that is incident to ordinary chords
of both colors with different other endpoints.  It excludes a common
pair and any other-color chord at a scalar branch triple.  In applying
the absolute-continuity clause, separate marginal absolute continuity
of the window increments is insufficient: the relative heights must
have the asserted joint law.  No statement is made at a singular center
where that joint local absolute continuity is unavailable.
\end{remark}
